% Einleitung — claude-tunnel bachelor thesis.
\chapter{Einleitung}
\label{ch:einleitung}

\section{Motivation}
\label{sec:motivation}

Wer heute einen Netzwerkdienst durch fremde Infrastruktur tunnelt -- sei es
ein \gls{ac:vpn}-Anbieter, ein Reverse-Proxy oder ein Relay-Netz --, vertraut
dem Betreiber dieser Infrastruktur weit mehr an, als technisch notwendig waere.
In den ueblichen Architekturen sieht der Vermittler mindestens die
Metadaten der Verbindung: welcher Client zu welchem Ziel-Hostnamen spricht, oft
auch, wie viel und wann. Bei TLS-terminierenden Proxies liegt zeitweise sogar
der Klartext vor. Dieses Wissen ist ein Angriffspunkt: es kann protokolliert,
verkauft, unter Zwang herausgegeben oder zur Zensur einzelner Ziele genutzt
werden. Anonymisierungsnetze wie Tor~\cite{dingledine2004tor} adressieren einen
Teil des Problems, verlagern das Vertrauen aber auf die Zusammensetzung des
Relay-Netzes und sind gegen aktive Netzsperren (etwa das Blockieren bekannter
Einstiegspunkte oder ganzer Transportprotokolle) nur begrenzt robust.

Die Gegenthese dieser Arbeit lautet: Ein Vermittler soll das, was er nicht
wissen muss, auch strukturell nicht wissen \emph{koennen}. Ein solcher Dienst
ist \gls{gl:providerblind} -- er relayt ausschliesslich Chiffretext, kennt weder
Ziel-Hostnamen noch Nutzlast und kann folglich weder das eine noch das andere
preisgeben. Vertrauen wird so nicht organisatorisch (durch Richtlinien und
Wohlverhalten des Betreibers) hergestellt, sondern kryptographisch erzwungen.

\section{Problemstellung}
\label{sec:problem}

Ein providerblinder Tunnel muss mehrere, teils gegenlaeufige Anforderungen
gleichzeitig erfuellen:

\begin{itemize}
  \item \textbf{Ende-zu-Ende-Vertraulichkeit gegen den Betreiber.} Die Nutzlast
    muss zwischen Client und Ziel (\emph{Origin}) verschluesselt sein, sodass die
    Vermittlungsebene nur opake Frames weiterreicht. Das schliesst eine
    TLS-Terminierung am Relay aus und verlangt eine eigene
    \gls{ac:e2e}-Schicht -- hier das Noise-Protokoll~\cite{perrin2018noise} ueber
    \gls{ac:quic}~\cite{rfc9000,rfc9001}.
  \item \textbf{Adressierung ohne Hostnamen.} Der Client muss den richtigen
    Endpunkt erreichen, ohne dem Betreiber einen Hostnamen zu nennen. Das
    verlangt eine Adressierung ueber opake Tokens und ein \gls{gl:rendezvous},
    das nicht auf Klarnamen beruht.
  \item \textbf{Missbrauchsresistenz ohne Identitaeten.} Gerade weil der Dienst
    blind ist, kann er Missbrauch nicht am Inhalt erkennen. Ein
    \gls{ac:pow}-Nachweis~\cite{back2002hashcash} verteuert die
    Verbindungsanbahnung und begrenzt so Flutangriffe, ohne Identitaeten zu
    fordern.
  \item \textbf{Zensur- und \gls{ac:nat}-Realitaeten.} Reale Netze blockieren
    zuweilen \gls{ac:udp} (und damit \gls{ac:quic}) oder liegen hinter
    \gls{ac:nat}. Der Tunnel braucht daher einen Transport-Fallback und einen
    direkten Pfad, der den Relay bei Erreichbarkeit umgeht~\cite{ford2005peer}.
\end{itemize}

Der Kern der Spannung liegt darin, diese Eigenschaften \emph{zusammen} zu
realisieren: Blindheit erschwert Missbrauchsabwehr, Ende-zu-Ende-Kryptographie
erschwert Vermittlung, und Zensurresistenz erschwert die Transportwahl.

\section{Zielsetzung und Forschungsfragen}
\label{sec:ziele}

Ziel der Arbeit ist der Entwurf, die prototypische Implementierung und die
empirische Evaluation eines providerblinden Tunnels, der die genannten
Anforderungen in einem lauffaehigen System vereint. Daraus leiten sich drei
Forschungsfragen ab:

\begin{description}
  \item[FF1 (Realisierbarkeit).] Laesst sich ein providerblinder Tunnel so
    konstruieren, dass der Betreiber nachweislich nur Chiffretext sieht, und
    zugleich Rendezvous, \gls{ac:pow}-Gating, direkter Pfad und
    Transport-Fallback funktionieren?
  \item[FF2 (Latenz-Overhead).] Welchen Roundtrip-Overhead verursacht der Tunnel
    gegenueber der reinen Netzlaufzeit, und wie skaliert er mit Verzoegerung und
    Paketverlust?
  \item[FF3 (Netz- und Betriebsartverhalten).] Wie verhalten sich die
    Betriebsarten One-shot, Streaming und Datagramm unter identischen
    Netzbedingungen, und welche Rolle spielt die \gls{ac:pow}-Schwierigkeit?
\end{description}

\section{Beitrag}
\label{sec:beitrag}

Die Arbeit liefert (i) einen in Rust implementierten Prototyp aus fuenf Crates
(gemeinsame Wire-Typen, Agent, Edge, Control-Plane, Client), der den vollen
Pfad Client\,$\rightarrow$\,Edge\,$\rightarrow$\,Agent\,$\rightarrow$\,Origin
mit \gls{ac:e2e}-Noise, \gls{ac:pow}-Rendezvous, direktem \gls{ac:p2p}-Pfad und
\gls{ac:tls}-ueber-TCP-Fallback traegt; (ii) ein rein containerbasiertes
Testbett, das die Topologie samt Stoergroessen (\texttt{tc netem}) emuliert; und
(iii) eine Parameterstudie am fertigen Produkt, die den Latenz-Overhead ueber
die drei Betriebsarten statistisch belastbar vermisst.

\section{Aufbau der Arbeit}
\label{sec:aufbau}

Kapitel~\ref{ch:einleitung} umreisst Motivation und Fragestellung.
Die weiteren Kapitel behandeln nacheinander die technischen Grundlagen
(Zero-Knowledge-Prinzip, \gls{ac:quic}, Noise, \gls{ac:pow}, \gls{ac:nat}),
verwandte Arbeiten, Anforderungen und Bedrohungsmodell, die Architektur, die
Implementierung, die Evaluation mit der Parameterstudie sowie Diskussion und
Ausblick.
